\section{Experiments}
\label{sec:experiments}

\subsection{Setup}

\paragraph{Model.} All variants use the \textsc{tiny} configuration from Table~\ref{tab:configs}: a 125M-parameter hybrid backbone with $\dmodel = 1024$, 12 layers (10 SSM, 2 attention), $\dstate = 16$, 8 query heads, and 1 shared KV head. All variants tie word embeddings and the LM head.

\paragraph{Data and tokenisation.} We pretrain on the FineWeb-Edu \texttt{sample-10BT} subset~\citep{lozhkov2024fineweb} tokenised with the GPT-2 BPE~\citep{radford2019language}, producing $\approx 10$B uint16 tokens stored as flat binary shards. Validation perplexity is computed on a held-out 50M-token slice of the same distribution.

\paragraph{Training.} We train with bfloat16 mixed precision, gradient checkpointing on the SSM scan, AdamW~\citep{loshchilov2017decoupled} in 8-bit~\citep{dettmers20218bit} with $\beta = (0.9, 0.95)$, weight decay $0.1$ (excluded from biases, layer norms, the SSM $A$-log and $D$ parameters, hormone magnitudes, and gates), peak learning rate $3 \times 10^{-4}$, $2{,}000$ warmup steps, and cosine decay to $10\%$ of peak. Sequence length is $1024$. Training uses DDP across 8 GPUs; full hyperparameters are listed in Appendix~\ref{sec:appendix}. Gradient clipping is applied at $1.0$.

\paragraph{Token budget.} Stage 1 (base pretraining) consumes $1$B tokens (\TBD{XXX} optimiser steps). For hormone routing variants, hormone extraction occurs at the $1$B-token milestone and Stage 2 contributes a further $1$B tokens, for a total of $2$B tokens across all variants. This isolates the effect of the modifications from the optimisation budget.

\paragraph{Hardware.} The fused-kernel variants (\textsc{Base}, \textsc{HR}, \textsc{HR-rand}, \textsc{HR-noext}, \textsc{HR-fixedgate}) are trained on $8\times$ NVIDIA RTX 4090/5090 (24\,GB) consumer GPUs. The reference-scan variants (\textsc{IA}, \textsc{Both}) are trained on $8\times$ NVIDIA RTX PRO 6000S (96\,GB) where the larger micro-batch (16 vs.\ 2) absorbs the per-step Python overhead of the reference scan. The total spend across all reported runs is approximately \TBD{\$XX}.

\subsection{Variants}

We compare seven variants holding the 2B-token optimisation budget fixed:
\begin{description}
    \item[\textsc{Base}] The unmodified hybrid backbone. No intracellular attention, no hormone routing.
    \item[\textsc{IA}] Hybrid backbone with intracellular attention enabled in every SSM layer (stride $k=32$); no hormone routing.
    \item[\textsc{HR}] Hybrid backbone with hormone routing enabled (hormones extracted at the $1$B-token milestone). No intracellular attention.
    \item[\textsc{Both}] Both intracellular attention and hormone routing enabled.
    \item[\textsc{HR-rand}] Hormone routing with the hormone bank $V$ replaced by random unit vectors of the same shape. Tests whether the specific extracted directions matter.
    \item[\textsc{HR-noext}] Hormone routing with $V$ extracted at step zero from a randomly initialised model. Tests whether extraction must follow Stage 1 training.
    \item[\textsc{HR-fixedgate}] Hormone routing with the output gate $g$ initialised to $1.0$ and detached from gradients. Tests the value of the zero-init learnable gate.
\end{description}

The intracellular attention variants (\textsc{IA} and \textsc{Both}) require the Python reference scan, which is substantially slower per training token than the fused kernel. To keep the optimisation budget fixed across all variants, these runs are conducted on a high-VRAM compute node ($8 \times 96$\,GB) where the larger micro-batch amortises the per-step Python overhead.

\subsection{Main results}

Table~\ref{tab:main} reports validation perplexity and downstream zero-shot accuracy on five standard benchmarks: HellaSwag~\citep{zellers2019hellaswag}, LAMBADA~\citep{paperno2016lambada}, ARC-Easy and ARC-Challenge~\citep{clark2018think}, and PIQA~\citep{bisk2020piqa}. We evaluate using direct log-likelihood scoring: for each multiple-choice example we compute $\log P(\text{continuation} \mid \text{context})$ for each candidate and predict the highest-scoring option. For LAMBADA we use greedy-argmax accuracy on the last word. This avoids any dependency on evaluation harness versions and allows exact reproduction from the model checkpoint alone.

\begin{table}[t]
\centering
\caption{Main results on \textsc{tiny} (125M parameters), 2B-token training budget. Lower perplexity is better; higher accuracy is better. HellaSwag reported as acc$_\text{norm}$ (length-normalised); ARC as acc$_\text{norm}$; PIQA as acc$_\text{norm}$; LAMBADA as exact-match accuracy. Best per column in \textbf{bold}.}
\label{tab:main}
\begin{tabular}{lcccccc}
\toprule
Variant & Val ppl $\downarrow$ & HellaSwag & LAMBADA & ARC-E & ARC-C & PIQA \\
\midrule
\textsc{Base}        & \TBD{X.XX} & \TBD{X.X} & \TBD{X.X} & \TBD{X.X} & \TBD{X.X} & \TBD{X.X} \\
\textsc{IA}          & \TBD{X.XX} & \TBD{X.X} & \TBD{X.X} & \TBD{X.X} & \TBD{X.X} & \TBD{X.X} \\
\textsc{HR}          & \TBD{X.XX} & \TBD{X.X} & \TBD{X.X} & \TBD{X.X} & \TBD{X.X} & \TBD{X.X} \\
\textsc{Both}        & \TBD{X.XX} & \TBD{X.X} & \TBD{X.X} & \TBD{X.X} & \TBD{X.X} & \TBD{X.X} \\
\midrule
\textsc{HR-rand}     & \TBD{X.XX} & \TBD{X.X} & \TBD{X.X} & \TBD{X.X} & \TBD{X.X} & \TBD{X.X} \\
\textsc{HR-noext}    & \TBD{X.XX} & \TBD{X.X} & \TBD{X.X} & \TBD{X.X} & \TBD{X.X} & \TBD{X.X} \\
\textsc{HR-fixedgate}& \TBD{X.XX} & \TBD{X.X} & \TBD{X.X} & \TBD{X.X} & \TBD{X.X} & \TBD{X.X} \\
\bottomrule
\end{tabular}
\end{table}

\TBD{Headline narrative: e.g.\ ``\textsc{Both} attains the lowest validation perplexity at $X.XX$, an improvement of $\Delta$ over \textsc{Base}. The two contributions compose: \textsc{IA} alone gives $\Delta_1$, \textsc{HR} alone gives $\Delta_2$, and \textsc{Both} gives approximately $\Delta_1 + \Delta_2$, suggesting the mechanisms operate independently. The hormone-routing ablations show [pattern across rand/noext/fixedgate].''}
